% ===== 原文 pp.130--134 =====
% [原文 p.130]
一个没有动力学引力的量子场论，我们必须取极限
\begin{equation}
  G_N \to 0 \iff \vol(M_D) \to \infty\,. \tag{6.3}
\end{equation}

\begin{figure}[htbp]
\centering
\includegraphics[width=0.62\textwidth]{fig28.png}
\figtitle{28}{把膜包裹在非紧闭链上（图中蓝色）给出非动力学缺陷；另一方面，把膜包裹在紧闭链上（图中青色）给出动力学激发。}
\end{figure}

\para{缺陷与动力学激发。}
我们要考察的问题是：在这个几何工程框架中如何构造高阶形式对称性。为此我们构造各种维数的未屏蔽算符构成的群\footnote{直到不久前，这一方向在量子场论的弦论实现范围内的研究，主要着眼于辨认广义对称性下（未必是拓扑的）带电缺陷的弦论起源。至于\emph{拓扑}对称性产生子的起源，[62,103] 中已有一些用通量算符表述的早期讨论。较新的工作 [136, 138, 139, 154] 说明，在适当的拓扑极限下，几何工程与全息构造中包裹的膜都描述拓扑对称性产生子。}。这有两类贡献：
\begin{itemize}[leftmargin=1.4em,itemsep=1pt,topsep=1pt]
\item 包裹内禀几何 $M_D$ 中\emph{非紧} $k$-闭链的 $p$-膜，给出 $M_d$ 上 QFT $\mathcal{T}$ 中的 $(p-k+1)$ 维缺陷。它们由 $M_D$ 相对于其边界 $\partial M_D$ 的相对同调 $H_k(M_D,\partial M_D)$ 分类。详见下面的方框。
\item 包裹\emph{紧} $k$-闭链的 $p$-膜，给出场论 $\mathcal{T}$ 中的 $(p-k+1)$ 维动力学激发。它们由同调群 $H_k(M_D)$ 分类。
\end{itemize}
把膜紧致化在非紧闭链上给出缺陷，是因为紧致化后的物体在有效 QFT $\mathcal{T}$ 中的质量由内禀闭链的体积决定，而质量无穷大的物体可以视为缺陷。图 28 说明了这一点。

% [Tutor Note: 两个 bullet 的分工是 §6.1 的核心：非紧闭链 $\leftrightarrow$ 缺陷（其质量 $\sim$
%  该闭链体积，随 $G_N\to0$ 时体积 $\to\infty$ 而变成无穷大质量，即非动力学），
%  紧闭链 $\leftrightarrow$ 有限质量的动力学激发。由此「未屏蔽算符」的群才由 (6.4) 的相对同调给出。]

% [原文 p.131]
\begin{notebox}{相对同调}
按定义，$H_k(M_D,\partial M_D)$ 由相对闭链模相对边界给出。相对闭链的定义放宽了通常闭链的定义：相对链 $x \in C_k(M_D)$ 是相对闭链，若 $\partial x \in C_{k-1}(\partial M_D)$。在相对同调中，相对闭链 $x$ 是平凡的，若 $x = \partial y + z$，其中 $y \in C_{k+1}(M_D)$、$z \in C_k(\partial M_D)$。
\end{notebox}

\para{屏蔽与 p-形式对称性。}
把缺陷被动力学激发屏蔽这件事考虑在内，我们发现包裹 $M_D$ 中 $k$-闭链的 $p$-膜给出如下未屏蔽的 $(p-k+1)$ 维算符构成的群
\begin{equation}
  \mathsf{D}_{p,k} = H_k(M_D,\partial M_D;\Z)/H_k(M_D;\Z)\,, \tag{6.4}
\end{equation}
把所有这些合起来，就得到所谓的\emph{缺陷群}（defect group）[62,63,201]
\begin{equation}
  \mathsf{D} = \bigoplus_{p,k} \mathsf{D}_{p,k} = \bigoplus_i \mathsf{D}_i\,, \tag{6.5}
\end{equation}
其中右端我们按维数为 $i$ 的缺陷重新组织了 $\mathsf{D}$。$\mathsf{D}$ 上一般有非退化配对，它是一个双同态
\begin{equation}
  \langle \cdot , \cdot \rangle : \mathsf{D} \times \mathsf{D} \to U(1)\,, \tag{6.6}
\end{equation}
它刻画 $\mathsf{D}$ 中缺陷之间的相互非局域性。这种相互非局域性反映了：并非所有这些缺陷都能同时处于高阶形式对称性的未扭曲扇区中。有些缺陷必须附接在产生高阶形式对称性的拓扑算符上，而当缺陷穿过这些拓扑算符、按其荷获得相位时，非局域性就出现了。

对于一致地选取哪些缺陷被扭曲、哪些不被扭曲，有多种选择。每种这样的选择对应于所谓的\textbf{极化}（polarization）或\textbf{拉格朗日子群}（Lagrangian subgroup），即
\begin{equation}
  \Lambda = \bigoplus_i \Lambda_i\,, \tag{6.7}
\end{equation}
其中每个 $\Lambda_i$ 是一个子群
\begin{equation}
  \Lambda_i \subseteq \mathsf{D}_i\,, \tag{6.8}
\end{equation}
使得 $\Lambda$ 是使该配对平凡化的极大子群
\begin{equation}
  \langle \cdot , \cdot \rangle|_\Lambda = 1\,. \tag{6.9}
\end{equation}

% [Tutor Note: (6.4) 里把 $H_k(M_D)\subset H_k(M_D,\partial M_D)$ 当作子群来取商，
%  这一步只在同调无挠时干净成立——这正是 p.133 脚注 8 所声明的假设
%  （「这里假设 $M_D$ 的同调中没有挠率」）。若 $M_D$ 有挠，需要用长正合列
%  (6.12) 的写法 $\mathsf{D}_{p,k}=\mathrm{im}(g_k)$ 才不致出错。]

% [原文 p.132]
不过应当注意，由于 't Hooft 反常的存在，并非所有这样的 $\Lambda$ 选择都被允许。每个被允许的 $\Lambda$ 选择给出所构造的 $d$ 维 QFT 的一个「全局形式」。人们可以通过对下面描述的高阶形式对称性作离散规范化（可能带离散挠率），在不同全局形式之间循环。

给定一个被允许的 $\Lambda$，所得的 $d$ 维 QFT 具有高阶形式对称性
\begin{equation}
  \mathcal{S}_\Lambda = \prod_i G^{(i)}\,, \tag{6.10}
\end{equation}
其中 $i$-形式对称群为
\begin{equation}
  G^{(i)} = \widehat{\Lambda}_i\,. \tag{6.11}
\end{equation}

\begin{notebox}{与 swampland 猜想的一致性}
按照 swampland 猜想之一，一个自洽的量子引力理论不能有任何整体对称性，包括高阶形式对称性。在几何工程的语境中，当内禀紧致化流形是紧的时候会出现量子引力理论，因为此时紧致化理论带有非零的 $G_N$。我们所讨论的所有高阶形式对称性都与内禀流形中的非紧闭链相关联，因此这里 swampland 猜想是满足的。

更一般地，关于 swampland 与高阶形式对称性交界处工作的一些样本，见 [208--214]。
\end{notebox}

\para{用边界来描述。}
我们可以用内禀流形的边界 $\partial M_D$，借助相对同调的长正合列来描述 $\mathsf{D}_{p-k+1}$
\begin{equation}
\begin{aligned}
\cdots \to H_i(\partial M_D) &\xrightarrow{h_i} H_i(M_D) \xrightarrow{f_i} H_i(M_D,\partial M_D)\\
&\xrightarrow{g_i} H_{i-1}(\partial M_D) \xrightarrow{h_{i-1}} H_{i-1}(M_D) \to \cdots
\end{aligned}\tag{6.12}
\end{equation}
这给出
\begin{equation}
  H_k(M_D,\partial M_D)/H_k(M_D) = \mathrm{im}(g_k) = \ker(h_{k-1})\,, \tag{6.13}
\end{equation}
从而
\begin{equation}
  \mathsf{D}_{p,k} = \ker(h_{k-1})\,. \tag{6.14}
\end{equation}
于是 $\mathsf{D}_{p,k}$ 有一个自然的解释：它是边界 $\partial M_D$ 上那些被推入体 $M_D$ 后变得平凡的 $(k-1)$-闭链。

% [Tutor Note: (6.12) 的箭头方向与标准配对长正合列一致：
%  $H_i(\partial M_D)\xrightarrow{h_i}H_i(M_D)$ 由包含映射诱导，
%  $H_i(M_D)\xrightarrow{f_i}H_i(M_D,\partial M_D)$ 由链群间的自然映射诱导，
%  $H_i(M_D,\partial M_D)\xrightarrow{g_i}H_{i-1}(\partial M_D)$ 是连接同态。
%  于是 $\mathrm{im}(g_k)=\ker(h_{k-1})$ 恰是正合性在 $H_{k-1}(\partial M_D)$ 处的表述，
%  这正是 (6.13)；用它可以检查 (6.14) 的下标：$\mathsf{D}_{p,k}$ 对应的是 $(k-1)$-闭链。]

\begin{examplebox}{例 6.1：7d $\mathcal{N}=1$ SYM}
考虑 M-理论在
\begin{equation}
  M_{\mathcal{D}=4} = \C^2/\Gamma_{ADE}\,, \tag{6.15}
\end{equation}
其中 $\Gamma_{ADE}$ 是 $\SU(2)$ 的有限子群，采用基于

% [原文 p.133]
McKay 对应的 ADE 分类。通过恰当地选取渐近通量 [62,63]，该构型工程化出规范代数为 $\mathfrak{g}_{ADE}$ 的纯 7d $\mathcal{N}=1$ 超对称 Yang--Mills 理论。相应的单连通群记为 $G_{ADE}$。由前面讨论的规范理论分析可知，这些理论具有缺陷群
\begin{equation}
  \mathsf{D} = \mathsf{D}_1 \oplus \mathsf{D}_4\,, \tag{6.16}
\end{equation}
使得
\begin{equation}
  \mathsf{D}_1 \cong \mathsf{D}_4 \cong Z(G_{ADE})\,, \tag{6.17}
\end{equation}
其中 $Z(G_{ADE})$ 是 $G_{ADE}$ 的中心。这里 $\mathsf{D}_1$ 是未屏蔽 Wilson 线算符构成的群，$\mathsf{D}_4$ 是未屏蔽 't Hooft 算符构成的群。

现在用几何工程分析来计算缺陷群。内禀流形的边界是
\begin{equation}
  \partial(\C^2/\Gamma_{ADE}) = S^3/\Gamma_{ADE}\,, \tag{6.18}
\end{equation}
其一阶同调群为
\begin{equation}
  H_1(S^3/\Gamma_{ADE}) = \mathrm{Ab}[\Gamma_{ADE}] = Z(G_{ADE})\,, \tag{6.19}
\end{equation}
其中 $\mathrm{Ab}[\Gamma_{ADE}]$ 是（一般为非阿贝尔的）群 $\Gamma_{ADE}$ 的交换化，它同构于 $Z(G_{ADE})$。所有这些 1-闭链推入体中都收缩。因此把 M2 与 M5 膜包裹在它们上面，就得到 (6.16) 中的贡献 $\mathsf{D}_1$ 与 $\mathsf{D}_4$。
\end{examplebox}

% [Tutor Note: (6.19) 可用 ADE 中的 $D_4$ 核对：$\Gamma_{D_4}=Q_8$（四元数群），
%  其交换化 $\mathrm{Ab}[Q_8]=\Z_2\times\Z_2$，而 $\Z(\mathrm{Spin}(8))=\Z_2\times\Z_2$，
%  二者一致；再看 $A_n$：$\Gamma=\Z_{n+1}$，$\mathrm{Ab}=\Z_{n+1}=\Z(\mathrm{SU}(n+1))$。
%  注意 (6.17) 说 $\mathsf{D}_1\cong\mathsf{D}_4$，这要求两个方向的 Wilson/'t Hooft 算符都未被屏蔽。]

\para{用相交数来描述。}
由 Poincaré 对偶，有\footnote{这里我们假设 $M_D$ 的同调中没有挠率。}
\begin{equation}
  H_k(M_D,\partial M_D) = H_{D-k}(M_D)\,, \tag{6.20}
\end{equation}
其中右端是内禀流形 $M_D$ 上紧 $(D-k)$-闭链的同调群。相应的商现在可以写为
\begin{equation}
  \mathsf{D}_{p,k} = H_k(M_D,\partial M_D)/H_k(M_D) = H_{D-k}(M_D)/H_k(M_D)\,, \tag{6.21}
\end{equation}
这要求一个单同态
\begin{equation}
  i : H_k(M_D) \to H_{D-k}(M_D)\,. \tag{6.22}
\end{equation}
映射 $i$ 把一个 $k$-闭链 $C_k$ 送到 $(D-k)$-闭链
\begin{equation}
  i(C_k) = \sum_j \left(C_k \cdot C^j_{D-k}\right) C^j_{D-k}\,, \tag{6.23}
\end{equation}
其中 $C^j_{D-k}$ 是 $H_{D-k}(M_D)$ 中 $(D-k)$-闭链的一组基，而 $C_k \cdot C^j_{D-k}$ 是 $C_k$ 与 $C^j_{D-k}$ 在 $M_D$ 中的相交数。

% [原文 p.134]
\begin{examplebox}{例 6.2：5d $\mathcal{N}=1$ SCFT：$E_1$ 理论}
一个有趣且著名的 5d $\mathcal{N}=1$ SCFT 是 $E_1$ 理论，它的特征性质是：在质量形变之后，它流向 5d $\mathcal{N}=1$ 纯 $\SU(2)$ 规范理论（离散 theta 角平凡）。

该理论有一个几何工程构造，其中内禀流形是含单个紧 4-闭链的非紧 Calabi--Yau 三流形 $M_6$，该 4-闭链的拓扑为
\begin{equation}
  C_4 = \mathbb{P}^1 \times \mathbb{P}^1\,. \tag{6.24}
\end{equation}
唯一的紧 2-闭链是 $C_4$ 内的两个 $\mathbb{P}^1$，二者的相交数都是
\begin{equation}
  \mathbb{P}^1 \cdot C_4 = -2\,. \tag{6.25}
\end{equation}
于是，对这个几何有
\begin{equation}
  H_2(M_6,\partial M_6)/H_2(M_6) = \Z_2\,. \tag{6.26}
\end{equation}
把 M2 膜包裹在这些非紧 2-闭链上，我们得到未屏蔽线算符构成的群
\begin{equation}
  \mathsf{D}_1 = \Z_2\,, \tag{6.27}
\end{equation}
这些线算符在 $E_1$ 理论中处于未扭曲扇区，于是我们推断出 $E_1$ 理论的 1-形式对称群
\begin{equation}
  G^{(1)} = \widehat{\mathsf{D}}_1 = \Z_2\,. \tag{6.28}
\end{equation}
\end{examplebox}

% [Tutor Note: (6.25) 的 $-2$ 是非紧 CY 三流形里紧 4-闭链 $C_4=\mathbb{P}^1\times\mathbb{P}^1$
%  上 $\mathbb{P}^1$ 的自相交数，与 $E_1$ 理论流向纯 $\SU(2)$ Yang--Mills 后
%  瞬子数/规范耦合的几何解释相容；可用 (6.26) 反推：两个紧 2-闭链贡献 $\Z\oplus\Z$，
%  再按 (6.26) 除以紧的 $H_2(M_6)$ 后只剩 $\Z_2$，说明这唯一的紧 4-闭链给出的正是 $\Z_2$ 的线算符群。]

\subsection{全息中的广义对称性}
本小节说明广义对称性如何出现在全息框架中\footnote{早期工作见 [95, 101]；较近的应用例如见 [99, 106, 133, 136, 138, 142, 161,175,177,194,215]。}。本小节的内容与 §4.5 关于对称性拓扑场论\footnote{关于全息语境中 SymTFT 与拓扑体项之间精确对应关系的深入讨论，例如见 [175]。}的内容密切相关。

\para{来自全息的对偶 SymTFT。}
考虑一个背景几何为
\begin{equation}
  \mathrm{AdS}_{d+1} \times X_{D-d-1}\,, \tag{6.29}
\end{equation}
的构造，其中 $D = 10, 11$，$X$ 表示某个内禀流形\footnote{这里限制到 AdS 解只是为了示意。例如 Klebanov--Strassler 解 [216] 就不是这种类型，但已用下文的技术在 [136,177,215] 中研究过。}。在把 10/11 维超引力作量纲约化、并只关注边界处的主导贡献之后，
